% TOP_PROBLEM: {"id": 30006682, "title": "Hardy-Littlewood Prime k-Tuple Conjecture", "category_id": 1, "category": {"id": 1, "name": "number_theory", "display_name": "Number Theory", "description": "Properties of integers, prime numbers, Diophantine equations.", "slug": "number-theory", "order_index": 1, "created_at": "2026-07-31T15:26:25.670Z"}, "status": "open"}
% FIRST_POSTED: null
% !TeX program = lualatex
% Compile twice: lualatex open_problems_500.tex
% All references and prose are embedded; no BibTeX database or external inputs.
\documentclass[11pt,a4paper]{article}
\usepackage[margin=24mm,headheight=14pt]{geometry}
\usepackage{fontspec}
\setmainfont{Latin Modern Roman}
\setsansfont{Latin Modern Sans}
\setmonofont{Latin Modern Mono}
\usepackage{amsmath,amssymb,mathtools}
\usepackage{unicode-math}
\setmathfont{Latin Modern Math}
\usepackage{microtype}
\usepackage{enumitem,xurl,needspace}
\usepackage[dvipsnames]{xcolor}
\usepackage{hyperref}
\hypersetup{unicode=true,colorlinks=true,linkcolor=MidnightBlue,urlcolor=MidnightBlue,
 pdftitle={Mathematical Research Notebook},pdfauthor={Source-annotated compilation},bookmarksnumbered=true}
\usepackage{bookmark}
\usepackage{tocloft}
\setlength{\cftsecnumwidth}{3em}
\cftsetpnumwidth{2.5em}
\cftsetrmarg{4em}
\usepackage{titlesec}
\titleformat{\section}{\Large\bfseries\raggedright}{\thesection}{0.7em}{}
\usepackage{fancyhdr}
\pagestyle{fancy}\fancyhf{}
\fancyhead[L]{\small\sffamily Open Problems Math | Research notebook}
\fancyhead[R]{\small 27 September 2026}
\fancyfoot[C]{\thepage}
\setlength{\parindent}{0pt}
\setlength{\parskip}{4pt plus 1pt}
\setlength{\emergencystretch}{3em}
\setcounter{tocdepth}{1}
\setcounter{secnumdepth}{1}
\setlist[itemize]{leftmargin=3em,itemsep=3pt,topsep=4pt,parsep=0pt}
\allowdisplaybreaks
\newcommand{\N}{\mathbb{N}}
\newcommand{\Z}{\mathbb{Z}}
\newcommand{\Q}{\mathbb{Q}}
\newcommand{\R}{\mathbb{R}}
\newcommand{\C}{\mathbb{C}}
\newcommand{\F}{\mathbb{F}}
\newcommand{\A}{\mathbb{A}}
\newcommand{\PP}{\mathbb P}
\newcommand{\E}{\mathbb{E}}
\newcommand{\eps}{\varepsilon}
\newcommand{\dd}{\,\mathrm d}
\newcommand{\sref}[2]{\hyperlink{source:#1:#2}{[S#2]}}
\newcommand{\eref}[2]{\hyperlink{extra:#1:#2}{[E#2]}}
% Turn the next switch off for a shorter reading copy. The quotations preserve
% every exactTarget field, including qualifications omitted from a short summary.
\newif\ifshowcatalogue
\showcataloguetrue
\newcommand{\cataloguescope}[1]{\ifshowcatalogue
 \paragraph{Exact scope in the supplied catalogue.}
 \begin{quote}\small #1\end{quote}\fi}
\usepackage{etoolbox}
\AtBeginEnvironment{itemize}{\small}
% Standalone entry; original section content is delimited for verification.
\begin{document}
\setcounter{section}{0}
%% BEGIN ORIGINAL SECTION CONTAINER
% ================================================================
% problemId: 30006682
% Canonical problem ID: 30006682
% exactTarget SHA-256: 32037bf72581c571bc62cac549ecd07c646c24fcc5e7b7fdf122154f45e36357
\clearpage
\section*{\texorpdfstring{Hardy-Littlewood Prime k-Tuple Conjecture}{Hardy-Littlewood Prime k-Tuple Conjecture}}\label{30006682}
\subsection{Definitions and mathematical statement}
For a finite set $H=\{h_1,\ldots,h_k\}\subset\Z$ of distinct offsets, let $\nu_p(H)$ be the number of residues it occupies modulo $p$. Admissibility means $\nu_p(H)<p$ for every prime $p$. Define
\[
 \mathfrak S(H)=\prod_p\frac{1-\nu_p(H)/p}{(1-1/p)^k},\qquad
 \pi_H(x)=\#\{1\le n\le x:n+h_i\text{ prime for every }i\}.
\]
For each fixed admissible $H$, the conjecture is $\pi_H(x)\sim\mathfrak S(H)x/(\log x)^k$ as $x\to\infty$. Here $f\sim g$ means $f/g\to1$; this is stronger than infinitude alone.
\par\smallskip\textit{Formulation and concept references:} \sref{30006682}{1}.
\cataloguescope{For every finite admissible set H=\{h\_1,...,h\_k\} of distinct integers, prove that the number of n <= x for which every n+h\_i is prime is asymptotic to S(H)x/(log x)\textasciicircum{}k, where S(H) is the Hardy-Littlewood singular series, or disprove this formula.}
\subsection{Short English statement}
Does every locally admissible fixed pattern of primes occur with exactly the frequency predicted by its divisibility corrections?
%% BEGIN RESEARCH INSERTION
\subsection{Research attempt}
\paragraph{Current frontier.}
\textit{Research check: 27 September 2026.}
Maynard's multidimensional sieve proves bounded intervals containing many
primes, not the asymptotic for each prescribed tuple. Green--Tao's
linear-equations framework excludes the parallel linear parts in the
single-variable forms $n+h_i$. \eref{30006682}{1}, \eref{30006682}{2}.
The 2026 short-interval uniformity theorem in the supplied sources gives
Hardy--Littlewood consequences with an additional short average; it does
not remove that average for every fixed $H$. \sref{30006682}{4}.
The primorial proof claim in [S3] is retained as an original source, not
accepted as a verified proof. The following argument tests the precise
local-to-global passage on which such an approach would have to improve.

\paragraph{Attack route.}
We compare three possible routes: direct simultaneous sieving, cancellation
in von Mangoldt correlations, and averaging over offsets followed by
removal of the average. They respectively require control of a growing
sieve modulus, a genuine fixed-shift correlation estimate, or a uniform
bound on exceptional offsets. We pursue the first two together, so that
one can see exactly which information survives from the local factors.

\paragraph{Attempt: the exact local calculation and its range.}
Assume first $k\geq1$ and fix $H$. Write
\[
 W(z)=\prod_{p\leq z}p,\qquad
 R_z=\{a\bmod W(z):\gcd(\prod_i(a+h_i),W(z))=1\}.
\]
The Chinese remainder theorem gives
\[
 |R_z|=\prod_{p\leq z}(p-\nu_p(H)),\qquad
 \#\{n\leq x:n\bmod W(z)\in R_z\}
 =x\prod_{p\leq z}(1-\nu_p(H)/p)+O(|R_z|).
 \tag{18.1}
\]
The error follows by counting each residue class as $x/W(z)+O(1)$;
its implied constant is absolute. No independence assumption is used.
Normalizing the main term by the density for $k$ individually rough
integers gives exactly
\[
 \mathfrak S_z(H)=
 \prod_{p\leq z}\frac{1-\nu_p(H)/p}{(1-1/p)^k}.
 \tag{18.2}
\]
For every sufficiently large prime the offsets are distinct modulo $p$,
so $\nu_p(H)=k$. Taylor expansion at zero yields
\[
 \log\frac{1-k/p}{(1-1/p)^k}
 =-\frac{k(k-1)}{2p^2}+O_k(p^{-3}).
\]
The tail converges absolutely; admissibility makes each of the finitely
many remaining factors positive. This proves existence and positivity of
the singular series, rather than just treating its product formally.

For fixed $z$, (18.1) is excellent as $x\to\infty$, but it counts composite
survivors too. To make roughness equivalent to primality for positive
$n+h_i\leq x+\max_i|h_i|$, one would sieve through its square root (apart
from the small primes themselves). Formula (18.1) supplies no useful
relative error at that growing modulus. The problem is not the CRT
identity: it is the unproved cancellation in its errors when the cutoff
is increased with $x$. Taking $x\to\infty$ first at each fixed $z$ does
not justify evaluating a diagonal $z=z(x)$.

\paragraph{Attempt: an exact weighted reduction.}
Extend $\Lambda(t)$ by zero for integers $t\leq0$, and put
\[
 C_H(x)=\sum_{1\leq n\leq x}\prod_{i=1}^k\Lambda(n+h_i).
 \tag{18.3}
\]
A precise sufficient statement is
$C_H(x)=\mathfrak S(H)x+o_H(x)$ for every fixed admissible $H$.
Here is the removal of both prime powers and logarithmic weights.
The number of $n\leq x$ for which some $n+h_i=p^a$ with $a\geq2$ is
$O_H(\sqrt{x}\log x)$, by summing the elementary bounds
$(x+O_H(1))^{1/a}$ over $2\leq a\leq\log_2(x+O_H(1))$.
Every term in (18.3) is $O_H((\log x)^k)$. Therefore the total contribution
from such powers is $o_H(x)$.

Let $A_H(t)$ count genuine prime tuples up to $t$ and let
$w(t)=\prod_i\log(t+h_i)$ for $t$ larger than a fixed constant $t_0(H)$.
The weighted count $B_H(x)=\int_{t_0}^xw(t)\,dA_H(t)$ consequently satisfies
$B_H(x)=\mathfrak S(H)x+o_H(x)$. For large $t$,
\[
 w(t)\sim(\log t)^k,\qquad
 \frac{w'(t)}{w(t)^2}=O_H\!\left(\frac1{t(\log t)^{k+1}}\right).
\]
Stieltjes integration by parts gives
\[
 A_H(x)=\frac{B_H(x)}{w(x)}
       +\int_{t_0}^x B_H(t)\frac{w'(t)}{w(t)^2}\,dt+O_H(1).
 \tag{18.4}
\]
Since $B_H(t)=O_H(t)$, the integral is
$O_H(x/(\log x)^{k+1})$; one can bound it by splitting at $\sqrt{x}$.
Thus (18.3)'s proposed estimate implies exactly the stated unweighted
asymptotic, with no extra uniformity in $H$ demanded. Conversely the
unweighted asymptotic and integration by parts recover (18.3), after the
same prime-power estimate. This is an equivalent analytic target, not a
new proof of it. Negative offsets affect only finitely many initial terms.
If the empty tuple is admitted, its assertion is simply $\lfloor x\rfloor
\sim x$ and needs no argument.

\paragraph{Attempt: what an averaged theorem cannot by itself remove.}
The obstacle can be made quantitative without pretending to build fake
primes. For any positive comparison constants $s_h$, let nonnegative
model counts satisfy $b_h(X)=s_hX$ except for one fixed $h_0$, where
$b_{h_0}(X)=0$. Then
\[
 \frac1H\sum_{h=1}^H
 \left|\frac{b_h(X)}{s_hX}-1\right|=\frac1H
 \quad(H\geq h_0).
 \tag{18.5}
\]
Perfectly good mean convergence with $H\to\infty$ coexists with a
permanently bad fixed offset. This countermodel invalidates removal of an
average solely by nonnegativity. It is not a model of the primes and
not a counterexample to Hardy--Littlewood.

One possible strengthening would be a uniform estimate
$|C_H(x)-\mathfrak S(H)x|\leq \epsilon(x)x$ for every admissible tuple in
a box whose radius tends to infinity, with $\epsilon(x)\to0$ and the
number of offsets fixed. Every fixed tuple would eventually be included.
But that strengthening already contains the fixed-shift difficulty; the
current averaged statements do not supply it. A different route would
prove (18.3) by combining the CRT factors with cancellation beyond the
sieve's parity obstruction, rather than identifying survivors with primes.

\paragraph{Stress test.}
The convergent singular series is only a local density correction. The
error in (18.1) is not uniform enough for primality detection. No argument
here promotes higher uniformity in a finite-complexity system to the
parallel-shift system without checking its hypotheses. The models in
(18.5) test an inference, not arithmetic reality. Proving (18.3) even for
$H=\{0,2\}$ would yield the twin-prime assertion of excluded UnsolvedMath problem 11;
that dependency is explicitly retained, although no separate file for
UnsolvedMath problem 11 is produced.

\paragraph{Outcome.}
\textbf{Outcome: Exploratory approach with unresolved gap.}
The exact CRT count and weighted/unweighted equivalence isolate the
missing fixed-tuple asymptotic (18.3). Neither the sieve error nor the
exceptional-offset error is controlled at the required scale. No novelty
is claimed for the classical ingredients.
%% END RESEARCH INSERTION
\subsection{Sources}
\begin{itemize}\raggedright
\item[\textnormal{[S1]}]\hypertarget{source:30006682:1}{} k-Tuple Conjecture (MathWorld).
\url{https://mathworld.wolfram.com/k-TupleConjecture.html}.
{\small\textit{Catalogue formulation source.}}
\item[\textnormal{[S2]}]\hypertarget{source:30006682:2}{} Correlations of almost primes — Institut Henri Poincaré seminar abstract.
\url{https://www.ihp.fr/en/node/3096}.
{\small\textit{Catalogue research/status reference; not a proof endorsement.}}
\item[\textnormal{[S3]}]\hypertarget{source:30006682:3}{} Proof of the Hardy-Littlewood K-tuple Conjecture in the Distribution of Numbers Coprime with the Primorial — Tim Samshuijzen, May 2026 revision.
\url{https://rxiv.org/pdf/2501.0157v4.pdf}.
{\small\textit{Catalogue research/status reference; not a proof endorsement.}}
\item[\textnormal{[S4]}]\hypertarget{source:30006682:4}{} Higher uniformity of arithmetic functions in short intervals II. Almost all intervals — Matomäki, Radziwiłł, Shao, Tao and Teräväinen (2026).
\url{https://link.springer.com/article/10.1007/s00222-026-01408-6}.
{\small\textit{Catalogue research/status reference; not a proof endorsement.}}
%% BEGIN RESEARCH REFERENCES
\item[\textnormal{[E1]}]\hypertarget{extra:30006682:1}{} James Maynard, \emph{Small gaps between primes}, arXiv:1311.4600 (2013), Theorem 1.1 and the sieve setup.
\url{https://arxiv.org/abs/1311.4600}.
{\small\textit{Added research source. Bounded clusters versus a prescribed simultaneous prime pattern. Consulted 27 September 2026.}}
\item[\textnormal{[E2]}]\hypertarget{extra:30006682:2}{} Ben Green and Terence Tao, \emph{Linear Equations in Primes}, arXiv:math/0606088, introduction and finite-complexity hypotheses.
\url{https://arxiv.org/abs/math/0606088}.
{\small\textit{Added research source. Scope of the linear-forms results; parallel shifts are not covered by the finite-complexity case. Consulted 27 September 2026.}}
%% END RESEARCH REFERENCES
\end{itemize}

%% END ORIGINAL SECTION CONTAINER
\end{document}
